
\section{Addendum to Experiments}
\label{sec:appExperiments}

\subsection{Implementation Details}
\label{sec:appImplementation}

We describe some of our implementation details below.
\begin{itemize}[leftmargin=0.0in]
\vspace{-0.05in}
\item[]
\textbf{Domain and Fidelity space:}
Given a problem with arbitrary domain $\Xcal$ and $\Zcal$, we mapped them to
$[0,1]^d$ and $[0,1]^p$ by appropriately linear transforming the coordinates.

\item[]
\textbf{Initialisation:}
Following recommendations in~\citet{brochu12bo}
all GP methods were initialised with uniform random
queries with $\capital/10$ capital, where $\capital$ is the total capital used in the
experiment.
For \gpucbs and \gpeis all queries were initialised at $\zhf$ whereas
for the multi-fidelity methods, the fidelities were picked at random from
the available fidelities.

\item[]
\textbf{GP Hyper-parameters:}
Except in the first two experiments of Fig.~\ref{fig:toy},
the GP hyper-parameters were learned after initialisation by maximising
the GP marginal likelihood~\citep{rasmussen06gps} and then updated every 25 iterations.
We use an SE kernel for both $\phix$ and $\phiz$ and
instead of using one bandwidth for the entire fidelity space and domain, we learn
a bandwidth for each dimension separately.
We learn the kernel scale, bandwidths and noise variance using marginal likelihood.
The mean of the GP is set to be the median of the observations.

\vspace{-0.05in}
\item[]
\textbf{Choice of $\betat$:} $\betat$, as specified in Theorem~\ref{thm:main} has unknown
constants and tends to be too conservative in practice~\citep{srinivas10gpbandits}.
Following the recommendations in~\citet{kandasamy15addBO} we  set it to be of
the correct ``order''; precisely, $\betat = 0.5d\log(2\ell t + 1)$. 
Here, $\ell$ is the effective $L_1$ diameter of $\Xcal$ and is computed by scaling
each dimension by the inverse of the bandwidth of the SE kernel for that dimension.

\vspace{-0.05in}
\item[]
\textbf{Maximising $\utilt$:}
We used the \directs algorithm~\citep{jones93direct}.

\item[]
\textbf{Fidelity selection:}
Since we only worked in low dimensional fidelity spaces, the set $\candfidelst$
was constructed in practice by obtaining a finely sampled grid of $\Zcal$ and then
filtering out those which satisfied the $3$ conditions in~\eqref{eqn:candfidelst}.
In the second condition of~\eqref{eqn:candfidelst}, the threshold $\gamma(z)$ can
be multiplied up to a constant factor, i.e $c\gamma(z)$ without affecting our theoretical
results.
In practice, we started with $c=1$ but we updated it every 20 iterations via the following
rule: if the algorithm has queried $\zhf$ more than $75\%$ of the time in the last $20$
iterations, we decrease it to $c/2$ and if it queried less than $25\%$ of the time
we increase it to $2c$.
But the $c$ value is always clipped inbetween $0.1$ and $20$.
In practice we observed that the value for $c$ usually stabilised around
$1$ and $8$ although in some experiments it shot up to $20$.
Changing $c$ this way  resulted in slightly better performance in practice.

% \textbf{Multiple Bandwidths in fidelity selection:}
% In theory we assumed a single bandwidth $\hz$ for the fidelity space,
% but in practice we used a bandwidth for each dimensions.
% However, the algorithm uses the bandwidth primarily via the kernel $\phiz$.
% Therefore, we only need to 

\end{itemize}


\subsection{Description of Synthetic Functions}
\label{sec:appSynthetic}

The following are the synthetic functions used in the paper.
\begin{itemize}[leftmargin=0.0in]

\item[]
\textbf{GP Samples:}
For the GP samples in the first two experiments of Figure~\ref{fig:toy} we used an
SE kernel with bandwidth $0.1$ for $\phix$.
For $\phiz$ we used bandwidths $1$ and $0.01$ for the first and second experiments
respectively.
The function was constructed by obtaining the GP function values on a $50\times 50$ grid
in the two dimensional $\Zcal\times\Xcal$ space and then interpolating for evaluations
in between via bivariate splines.
For both experiments we used $\eta^2 = 0.05$ and the cost function $\cost(z) = 0.2 +
6z^2$.

\item[]
\textbf{Currin exponential function:}
The domain is the two dimensional unit cube $\Xcal = [0,1]^2$ and the fidelity 
was $\Zcal = [0, 1]$ with $\zhf = 1$.
We used $\cost(z) = 0.1 + z^2$, $\eta^2=0.5$ and,
\begin{align*}
\gunc(z, x) &= \left(1-0.1(1-z)\exp\left(\frac{-1}{2x_2}\right)\right)
  \left(\frac{2300x_1^3 + 1900x_1^2 + 2092x_1 + 60}{100x_1^3 + 
  500x_1^2 + 4x_1 + 20}\right).
\end{align*}

\item[]
\textbf{Hartmann functions:}
We used 
$g(z,x) = \sum_{i=1}^4 (\alpha_i - \alpha'_i(z)) \exp\big( -\sum_{j=1}^3 A_{ij}
(x_j-P_{ij})^2\big)$.
Here $A, P$ are given below for the $3$ and $6$ dimensional cases and
$\alpha = [1.0, 1.2, 3.0, 3.2]$.
Then $\alpha'_i$ was set as $\alpha'_i(z) = 0.1(1 - z_i)$ if $i\leq p$ for
$i=1,2,3,4$.
We constructed the $p=4$ and $p=2$ Hartmann functions for the $3$ and $6$ dimensional
cases respectively this way.
When $z = \zhf = \one_p$, this reduces to the usual Hartmann function commonly used
as a benchmark in global optimisation.
% The $M$\ssth fidelity function is
% $\funcM(x) = \sum_{i=1}^4 \alpha_i \exp\big( -\sum_{j=1}^3 A_{ij}
% (x_j-P_{ij})^2\big)$ where $A, P\in\RR^{4\times 3}$ are fixed matrices given below and
% $\alpha = [1.0, 1.2, 3.0, 3.2]$. For the lower fidelities we use the same form except
% change $\alpha$ to $\alpha^{(m)} = \alpha + (M-m)\delta$ where $\delta = [0.01, -0.01,
% -0.1, 0.1]$ and $M=3$. The domain is $\Xcal = [0,1]^3$.

For the $3$ dimensional case we used $\cost(z) = 0.05 + (1 - 0.05)z_1^3z_2^2$,
$\eta^2=0.01$ and,
\[
A = 
\begin{bmatrix}
3 & 10 & 30 \\
0.1 & 10 & 35 \\
3 & 10 & 30 \\
0.1 & 10 & 35 
\end{bmatrix},
\quad
P = 10^{-4} \times
\begin{bmatrix}
3689 & 1170 & 2673 \\
4699 & 4387 & 7470 \\
1091 & 8732 & 5547 \\
381 & 5743 & 8828
\end{bmatrix}.
\]
For the $3$ dimensional case we used $\cost(z) = 0.05 + (1 -
0.05)z_1^3z_2^2z_3^{1.5}z_4^{1}$, $\eta^2=0.05$ and,
\[
A = 
\begin{bmatrix}
10 & 3 & 17 & 3.5 & 1.7 & 8 \\
0.05 & 10 & 17 & 0.1 & 8 & 14 \\
3 & 3.5 & 1.7 & 10 & 17 & 8 \\
17 & 8 & 0.05 & 10 & 0.1 & 14 
\end{bmatrix},
\;\;
P = 10^{-4} \times
\begin{bmatrix}
1312 & 1696 & 5569 &  124 & 8283 & 5886 \\
2329 & 4135 & 8307 & 3736 & 1004 & 9991 \\
2348 & 1451 & 3522 & 2883 & 3047 & 6650 \\
4047 & 8828 & 8732 & 5743 & 1091 &  381 \\
\end{bmatrix}.
\]




\item[]
\textbf{Borehole function:}
This function was taken from~\citep{xiong13highAccuracy}.
We first let,
\begin{align*}
\func_2(x) &= \frac{2\pi x_3(x_4-x_6)}
  {\log(x_2/x_1) \left(1 + \frac{2x_7x_3}{\log(x_2/x_1)x_1^2x_8} +
    \frac{x_3}{x_5} \right)}, \\
\func_1(x) &= \frac{5 x_3(x_4-x_6)}
  {\log(x_2/x_1) \left(1.5 + \frac{2x_7x_3}{\log(x_2/x_1)x_1^2x_8} +
    \frac{x_3}{x_5} \right)}. 
\end{align*}
Then we define $g(z,x) = z\func_2(x) + (1-z)\func_1(x)$.
The domain of the function is $\Xcal = [0.05, 0.15; 100, 50K; 63.07K, 115.6K;$
$990, 1110; 63.1, 116; 700, 820;$ $1120, 1680; 9855, 12045]$
and $\Zcal=[0,1]$ with $\zhf = 1$.
We used $\cost(z) = 0.1 + z^{1.5}$ for the cost function and
$\eta^2=5$ for the noise variance.
% We first linear transform the variables to lie in $[0,1]^8$.

\item[]
\textbf{Branin function:}
We use the following function where $\Xcal=[[-5, 10], [0, 15]]^2$ and $\Zcal=[0,1]^3$.
\[
g(z,x) = a(x_2 - b(z_1)x_1^2 + c(z_2)x_1 - r)^2 + s(1-t(z))cos(x_1) + s,
\]
where $a = 1$, $b(z_1)=5.1/(4\pi^2) - 0.01(1-z_1)$
$c(z_2) = 5/\pi - 0.1(1-z_2)$, $r=6$, $s=10$ and $t(z_3)=1/(8\pi) + 0.05(1-z_3)$.
At $z=\zhf = \one_p$, this becomes the standard Branin function used as a benchmark in
global optimisation.
We used $\cost(z) = 0.05 + z_1^{3}z_2^{2}z_3^{1.5}$ for the cost function and
$\eta^2=0.05$ for the noise variance.


\end{itemize}



% \subsection{Details on the Real Experiments}

